\section{The Positron And The Tilted Path}

Once the electron has been read over the flat baseline, the opposite sign
cannot be obtained by wish. The grammar has already forbidden it there. A
positron reading requires a different comparison, not a second arbitrary
label. The same pair must be carried against a tilted baseline. Only then
does the charged residue face the reader with the opposite orientation, and
only then is the sign \(+1\) admissible.

Baseline dependence is not a weakness in the charge reading. It is what
makes the sign meaningful. A sign is not a magnitude; it is a relation to an
oriented reference. The flat path supplies one reference. The tilted path
supplies another. If the grammar erased the difference between them, it would
erase the very orientation it is trying to read. A positron is therefore not
the electron plus a new substance. It is the same pair grammar seen through a
comparison whose baseline has been tilted.

The tilt is a constrained act. It does not let the reader reverse any sign at
will. A permissible tilt must preserve the underlying pair, the second
variation, the boundary commitments, and the distinction between separated
parts and coupled residue. What changes is the orientation of the baseline
against which the residue is read. Under that changed comparison, the same
discriminated activity that read \(-1\) over the flat path reads \(+1\). The
grammar identifies the sign flip as baseline-relative, not as a second
independent event.

This is why the native baseline carries an asymmetry. Over the flat path the
electron is present and the positron is not. The positron appears only when
the reading has been tilted enough to make the opposite orientation
admissible. The asymmetry is not imposed by a story about preference. It is
forced by the grammar of what the baseline can carry. Native flat reading
permits the negative charge and forbids the positive one. Tilted reading
permits the positive charge by changing the orientation of comparison.

Loop language helps, provided it remains an illustration. An open path can
carry local comparisons and still return no charge, because it has not
closed the orientation into a recoverable holonomy. A closed loop may return
with a signed residue, even when local segments looked balanced. The
important point is not that every positron must be pictured as a loop. The
point is grammatical: sign appears when transport around the relevant
comparison closes with an orientation that the reader cannot quotient away.
The tilted path supplies that closure for the \(+1\) reading.

The Aharonov--Bohm example shows the same form at a later, physical register.
Two branches can travel through regions where the local field reading is
silent, then recombine with a phase difference because the connection carried
memory around the loop. The present chapter uses the example only to sharpen
the grammar: local flatness does not guarantee global triviality, and the
sign of a transported distinction can depend on the path by which the
baseline is closed. The positron is read when the path carries the opposite
orientation.

Annihilation gives a second illustration. A coincidence detector accepts two
opposed photon directions inside a finite window as evidence that opposite
charges have cancelled. That finite coincidence logic presupposes two signs:
the electron's \(-1\) and the positron's \(+1\). The detector does not create
the signs. It reads the cancellation that those signs permit. In the grammar
of this chapter, the positron sign has already been forced by the tilted
baseline before any later coincidence story can use it.

The crucial discipline is that the positron reading is neither arbitrary nor
available everywhere. If a proposed reading announces \(+1\) over the flat
path, it violates the baseline. If it announces \(+1\) after a permissible
tilt, it is not inventing a new charge; it is reading the same
second-variation activity under the opposite orientation. The grammar
therefore preserves both facts: the electron is native to the flat baseline,
and the positron is the tilted-baseline charge.

The pair production grammar now has a real split in sign. The signs cancel
when read together, but they are not indistinguishable before cancellation.
The flat baseline reads the electron. The tilted baseline reads the
positron. A reader that cannot carry the tilt cannot read the positron as a
physical record; it must treat the positive sign as unavailable or
metaphysical relative to that class. A reader that can carry the tilt reads
the opposite orientation without changing the underlying invariant.

Thus the positron is not a second foundation. It is the charged reading
forced by a change in baseline. The same grammar that forbids \(+1\) over
the flat path permits \(+1\) over the tilted path. Baseline dependence is the
condition of the reading, and the sign flip is the grammar's first clear
lesson that orientation belongs to the relation between invariant and
reader.

